\documentclass{subfiles}
\begin{document}\label{sec:arithmeticEncoding}
    Huffman encoding has a limitation: the codewords can have length less then one bit.
        This limitation of the Huffman encoding is evident when one of the source symbols
        have a higher probability then the others. 

    Proposed by Peter Elias, another of Fano's students, Arithmetic encoding
        output is no longer a string of codewords, 
        rather a unique encoding of the whole text.
        The core idea is that of encoding the whole as floating-point number 
        in the range [0, 1].

    The algorithm is relatively simple: initialize the range, split it into \(n\)
        subintervals each proportional to the probability of a symbol, 
        and after reading a symbol select the associated range and 
        update the subintervals accordingly. After the whole text has been read,
        output a value within it.
        Below we provide the pseudo-code of the procedure to encode a string using Arithmetic encoding.

    \subfile{../../TikZ/Procedure - Arithmetic encoding}

    \begin{example*}
        Consider the string \(S = abac\), how do we encode it? 
            By apply the encoding procedure we get the following.
            
            \subfile{../../TikZ/Figure - Example of Arithmetic encoding}
    \end{example*}


    The decoding process is symmetric: the range [0, 1] is divided in \(n\)
        subintervals, at each iteration the proper range is selected and 
        the intervals recomputed accordingly. The process is valid 
        because the procedure is syncronized; that is,
        both encoder and decoder know the source alphabet and the associated probabilities.
        Below the pseudo-code code of the procedure to decode a string encoded with Arithmetic encoding.

    \subfile{../../TikZ/Procedure - Arithmetic decoding.tex}
    \begin{example*}
        Consider \(x = \sfrac{39}{128}\). If the source symbols and probabilities 
            are that of the previous example, what is the decoded string?
            
            \subfile{../../TikZ/Figure - Example of Arithmetic decoding}
    \end{example*}
\end{document}
